\color{red}
\section{Równoległe i hiperrównoległe} % lang-pl
\label{sekcja_hiperboliczna}

XI

\color{red}

\begin{definition}
[geometria hiperboliczna] % lang-pl
	Geometria płaska oparta na aksjomatach geometrii neutralnej i na zaprzeczeniu aksjomatu Playfaira: przez punkt leżący poza prostą przechodzi więcej niż jedna prosta do niej równoległa. % lang-pl
\index{geometria!hiperboliczna} % lang-pl
\end{definition}

Nazwę \emph{hiperboliczna} wprowadzi Felix Klein w 1871 roku (zobacz sekcję \ref{sekcja_beltrami_cayley_klein}). % lang-pl
O geometrii hiperbolicznej wspomina się w $\Delta_{24}^7$. % lang-pl

Zakładamy, że wszystko, co wynika z aksjomatów bez piątego postulatu -- w tym pierwsze dwadzieścia osiem twierdzeń pierwszej księgi Euklidesa -- pozostaje w mocy \cite[s. 288, 291]{eves1_1972}. % lang-pl

Półproste $PX$ i $PY$ nazwiemy \emph{równoległymi} do $AB$ przez $P$, a proste przez $P$ leżące poza kątem $XPY$ zawierającym $PQ$ -- \emph{hiperrównoległymi} do $AB$ \cite[s. 291--292]{eves1_1972}. % lang-pl
\index{proste!hiperrównoległe} % lang-pl

Tę samą myśl -- istnienie równoległej granicznej -- Hilbert sformułuje jako osobny aksjomat: zaksjomatyzuje geometrię hiperboliczną bez liczb rzeczywistych i bez aksjomatu Archimedesa, dodając do trzynastu aksjomatów płaszczyzny Hilberta aksjomat równoległej granicznej -- dla prostej $l$ i punktu $P$ poza nią istnieje półprosta graniczna wychodząca z $P$, tworząca kąt ostry z prostopadłą $PQ$ do $l$. % lang-pl
Uczyni w ten sposób jawnym założenie, które milcząco przyjmowali Saccheri, Gauss, Bolyai i Łobaczewski; podkreśli też, że jego podejście, w przeciwieństwie do podejść Kleina i Bolyaia--Łobaczewskiego, nie sięga do trzeciego wymiaru \cite[s. 209]{greenberg_2010}. % lang-pl
% Źródła: Eves s. 295--298 (§7.3).

\color{black}